\documentclass[11pt]{article}

\usepackage[margin=1in]{geometry}
\usepackage{amsmath,amssymb,amsthm}
\usepackage{hyperref}

\newtheorem{theorem}{Theorem}
\newtheorem{corollary}{Corollary}
\theoremstyle{remark}
\newtheorem{remark}{Remark}

\newcommand{\C}{\mathbb C}
\newcommand{\N}{\mathbb N}
\newcommand{\NA}{\operatorname{NA}}
\newcommand{\Nuc}{\mathcal N}
\newcommand{\oplusone}{\oplus_1}

\title{A Hilbert-Pair Subcase For Norm-Attaining Projective Tensors}
\author{Run \texttt{fa\_banach\_001}, agent \texttt{agent\_lane\_06}}
\date{June 16, 2026}

\begin{document}
\maketitle

\begin{abstract}
Dantas, Jung, Roldan, and Rueda Zoca ask whether, for reflexive Banach spaces
\(X\) and \(Y\), norm-attaining tensors are dense in
\(X\widehat\otimes_\pi Y\).  They prove that every tensor in
\(H\widehat\otimes_\pi H\) attains its projective norm when \(H\) is a complex
Hilbert space.  We record the same conclusion for arbitrary pairs of complex
Hilbert spaces \(H_1,H_2\).  Consequently, every nuclear operator between two
complex Hilbert spaces attains its nuclear norm.  The result also extends to
finite \(\ell_1\)-sums of Hilbert spaces in the first variable and a Hilbert
space in the second variable.  This is a positive subcase, not a solution of
the general reflexive problem.
\end{abstract}

\section{Source Question}

The source paper introduces norm-attainment for projective tensors and nuclear
operators.  In Section 6, Question 6.1 asks:
\begin{quote}
Let \(X,Y\) be reflexive Banach spaces. Is it true that the set of all
norm-attaining tensors is dense in \(X\widehat\otimes_\pi Y\)?
\end{quote}

The same section recalls two positive cases already proved in the paper:
every tensor in \(H\widehat\otimes_\pi H\) attains its projective norm for a
complex Hilbert space \(H\), and norm-attaining tensors are dense in
\(L_p(\mu)\widehat\otimes_\pi L_q(\nu)\) for \(1<p,q<\infty\).

\section{Definitions}

Let \(X\) and \(Y\) be Banach spaces.  A tensor
\[
  z=\sum_{n=1}^\infty \lambda_n x_n\otimes y_n,
  \qquad \lambda_n>0,\quad x_n\in S_X,\quad y_n\in S_Y,
\]
is said to attain its projective norm if
\[
  \|z\|_\pi=\sum_{n=1}^\infty \lambda_n.
\]
Equivalently, the projective norm of \(z\) is realized by an absolutely
summable representation by simple tensors.

\section{Idea Of The Proof}

For Hilbert spaces, projective tensors are the same object as trace-class
operators.  A nuclear operator \(T:H_1\to H_2\) has a singular-value
decomposition
\[
  T=\sum_n s_n \langle\cdot,u_n\rangle v_n,
\]
where \((u_n)\) and \((v_n)\) are orthonormal systems and
\(\sum_n s_n\) is the trace-class, hence nuclear, norm.  Reading the same
formula as a tensor representation gives an optimal projective representation.

For finite \(\ell_1\)-sums, the projective tensor product distributes over the
\(\ell_1\)-sum.  Thus the norm of a tensor is the sum of the norms of its
Hilbert components, and concatenating the singular-value representations of
the components gives a single optimal representation.

\section{Hilbert Pairs}

\begin{theorem}\label{thm:hilbert-pair}
Let \(H_1\) and \(H_2\) be complex Hilbert spaces.  Then every tensor in
\(H_1\widehat\otimes_\pi H_2\) attains its projective norm.
\end{theorem}

\begin{proof}
Since Hilbert spaces have the approximation property, the canonical map
\[
  H_1\widehat\otimes_\pi H_2 \longrightarrow
  \Nuc(H_1^*,H_2)
\]
is an isometric isomorphism, where \(h\otimes k\) acts on \(H_1^*\) by
\(\varphi\mapsto \varphi(h)k\).  Thus it is enough to prove that every nuclear
operator \(T:H_1^*\to H_2\) attains its nuclear norm.

Let \(T:H_1^*\to H_2\) be nuclear.  Since \(H_1^*\) is again a Hilbert space,
nuclear operators are exactly trace-class operators.  Hence \(T\) has a
singular-value decomposition
\[
  T\varphi=\sum_{n=1}^{N} s_n \psi_n(\varphi) v_n,
\]
where \(N\in\N\cup\{\infty\}\), \(s_n>0\), \((\psi_n)\) is an orthonormal
system in \((H_1^*)^*\), \((v_n)\) is an orthonormal system in \(H_2\), and
\[
  \|T\|_{\Nuc}=\sum_{n=1}^{N}s_n .
\]
By reflexivity of \(H_1\), each \(\psi_n\in (H_1^*)^*\) is evaluation at a
unique unit vector \(u_n\in H_1\).  Under the tensor--nuclear identification,
this is precisely the tensor representation
\[
  z_T=\sum_{n=1}^{N} s_n\, u_n\otimes v_n,
\]
with all vectors of norm one and with coefficient sum equal to the projective
norm.  Therefore \(z_T\), and hence every element of
\(H_1\widehat\otimes_\pi H_2\), attains its projective norm.
\end{proof}

\begin{corollary}\label{cor:hilbert-nuclear}
Let \(H_1,H_2\) be complex Hilbert spaces.  Then every nuclear operator
\(T:H_1\to H_2\) attains its nuclear norm.
\end{corollary}

\begin{proof}
This is the operator form of the proof of Theorem \ref{thm:hilbert-pair}.
\end{proof}

\section[Finite ell-1 Sums Of Hilbert Spaces]{Finite \(\ell_1\)-Sums Of Hilbert Spaces}

\begin{theorem}\label{thm:l1-sum-hilbert}
Let \(H_1,\ldots,H_m\) and \(K\) be complex Hilbert spaces, and put
\[
  X=H_1\oplusone \cdots \oplusone H_m .
\]
Then every tensor in \(X\widehat\otimes_\pi K\) attains its projective norm.
\end{theorem}

\begin{proof}
The projective tensor product commutes isometrically with finite
\(\ell_1\)-sums:
\[
  X\widehat\otimes_\pi K
  \cong
  (H_1\widehat\otimes_\pi K)\oplusone\cdots\oplusone
  (H_m\widehat\otimes_\pi K).
\]
Indeed, the dual of the left-hand side is
\(\mathcal L(X,K^*)\), which is isometric to
\(\mathcal L(H_1,K^*)\oplus_\infty\cdots\oplus_\infty
\mathcal L(H_m,K^*)\), the dual of the right-hand side.

Let \(z\in X\widehat\otimes_\pi K\), and write
\[
  z=(z_1,\ldots,z_m),\qquad z_j\in H_j\widehat\otimes_\pi K.
\]
By Theorem \ref{thm:hilbert-pair}, each nonzero \(z_j\) admits a representation
\[
  z_j=\sum_{n=1}^{N_j} s_{j,n} h_{j,n}\otimes k_{j,n},
  \qquad
  \sum_{n=1}^{N_j}s_{j,n}=\|z_j\|_\pi,
\]
with \(s_{j,n}>0\), \(h_{j,n}\in S_{H_j}\), and \(k_{j,n}\in S_K\).
Embedding \(H_j\) into the \(j\)-th coordinate of \(X\) and concatenating these
representations gives
\[
  z=\sum_{j=1}^{m}\sum_{n=1}^{N_j}
  s_{j,n}\, \iota_j(h_{j,n})\otimes k_{j,n}.
\]
The coefficient sum is
\[
  \sum_{j=1}^{m}\sum_{n=1}^{N_j}s_{j,n}
  =\sum_{j=1}^{m}\|z_j\|_\pi
  =\|z\|_\pi,
\]
where the last equality is the \(\ell_1\)-sum norm.  Hence \(z\) attains its
projective norm.
\end{proof}

\section{Scope And Limitations}

Finite \(\ell_1\)-sums of Hilbert spaces are reflexive, so Theorem
\ref{thm:l1-sum-hilbert} is a positive subcase of Question 6.1.  The result is
still narrow: it does not address arbitrary reflexive spaces, and it gives no
counterexample or full criterion for density of norm-attaining tensors.

The argument also does not solve the source paper's Question 6.2 about whether
there are Banach spaces \(X,Y\) for which norm-attaining nuclear operators are
not dense in \(\Nuc(X,Y)\), nor Question 6.3 concerning property \(\alpha\) or
quasi-\(\alpha\).

\section{Verification Notes}

The proof uses only standard Hilbert-space trace-class theory and the standard
isometric identification between projective tensors and nuclear operators when
the relevant Hilbert space has the approximation property.  The source paper's
Proposition 3.8 is the diagonal case \(H_1=H_2\); Theorem
\ref{thm:hilbert-pair} is the same singular-value proof without requiring the
domain and range Hilbert spaces to coincide.

Cheap run indexes were searched for prior packets on arXiv:2006.09871 and for
the relevant norm-attaining tensor/nuclear-operator keywords.  Bounded web
searches found the source paper but no later exact answer to Question 6.1.
Novelty confidence should be considered modest because this is an elementary
extension of a standard trace-class argument.

\section{References}

\begin{itemize}
  \item S. Dantas, M. Jung, O. Roldan, and A. Rueda Zoca,
  \emph{Norm-attaining tensors and nuclear operators}, arXiv:2006.09871.
  \item I. Gohberg, S. Goldberg, and M. A. Kaashoek,
  \emph{Classes of Linear Operators}, Vol. I, cited in the source paper for
  the singular-value representation and trace-class norm formula.
  \item R. A. Ryan, \emph{Introduction to Tensor Products of Banach Spaces},
  cited in the source paper for the tensor/nuclear identifications and
  projective tensor basics.
\end{itemize}

\end{document}
